% GRA on Megillas Esther — Complete English Rendition
% XeLaTeX document with bilingual Hebrew/English
\documentclass[11pt,twoside,openany]{book}

% === Geometry ===
\usepackage[
  papersize={7in,10in},
  inner=0.85in, outer=0.65in,
  top=0.75in, bottom=0.75in,
  headheight=14pt
]{geometry}

% === Fonts ===
\usepackage{fontspec}
\usepackage{polyglossia}
\setdefaultlanguage{english}
\setotherlanguage{hebrew}

% Main English font
\setmainfont{EB Garamond}[
  Numbers=OldStyle,
  Ligatures=TeX
]
\setsansfont{Libertinus Sans}

% Hebrew font
\newfontfamily\hebrewfont{Frank Ruehl CLM}[
  Script=Hebrew,
  Scale=1.1
]
% For pasuk text (larger, with taamim)
\newfontfamily\pasukfont{Ezra SIL}[
  Script=Hebrew,
  Scale=1.2
]

% === Headers & Footers ===
\usepackage{fancyhdr}
\pagestyle{fancy}
\fancyhf{}
\fancyhead[LE]{\small\textsc{the vilna gaon on megillas esther}}
\fancyhead[RO]{\small\textit{\leftmark}}
\fancyfoot[CE,CO]{\thepage}
\renewcommand{\headrulewidth}{0.4pt}

% === Section formatting ===
\usepackage{titlesec}

% Chapter style
\titleformat{\chapter}[display]
  {\normalfont\huge\bfseries\centering}
  {}
  {0pt}
  {}
\titlespacing*{\chapter}{0pt}{-20pt}{20pt}

% Section = verse grouping
\titleformat{\section}[hang]
  {\normalfont\Large\bfseries}
  {}
  {0pt}
  {}
\titlespacing*{\section}{0pt}{18pt}{8pt}

% Subsection = peshat/remez headers
\titleformat{\subsection}[hang]
  {\normalfont\large\scshape}
  {}
  {0pt}
  {}
\titlespacing*{\subsection}{0pt}{12pt}{4pt}

% === Custom environments ===
\usepackage{mdframed}
\usepackage{xcolor}

% Color palette
\definecolor{pasukbg}{HTML}{F5F0E8}
\definecolor{pasukframe}{HTML}{8B7355}
\definecolor{peshatcolor}{HTML}{1a3a5c}
\definecolor{remezcolor}{HTML}{5c1a3a}
\definecolor{notecolor}{HTML}{3a5c1a}
\definecolor{lightgray}{HTML}{E8E8E8}

% Pasuk box
\newmdenv[
  backgroundcolor=pasukbg,
  linecolor=pasukframe,
  linewidth=1pt,
  roundcorner=3pt,
  innertopmargin=8pt,
  innerbottommargin=8pt,
  innerleftmargin=12pt,
  innerrightmargin=12pt,
  skipabove=6pt,
  skipbelow=6pt
]{pasukbox}

% Commentary note box (for באר אברהם / רד"ל)
\newmdenv[
  backgroundcolor=white,
  linecolor=lightgray,
  linewidth=0.5pt,
  topline=true,
  bottomline=false,
  leftline=false,
  rightline=false,
  innertopmargin=6pt,
  innerbottommargin=6pt,
  skipabove=8pt,
  skipbelow=4pt
]{notebox}

% === Macros ===
% Pasuk (verse text)
\newcommand{\pasuk}[2]{%
  \begin{pasukbox}
  \begin{center}
  {\pasukfont\large\begin{hebrew}#1\end{hebrew}}
  \par\vspace{4pt}
  {\small\itshape #2}
  \end{center}
  \end{pasukbox}
}

% Verse reference
\newcommand{\vref}[1]{\textbf{#1}\quad}

% Hebrew inline
\newcommand{\heb}[1]{{\hebrewfont\begin{hebrew}#1\end{hebrew}}}

% Hebrew term in English context
\newcommand{\hebterm}[2]{\textit{#1} (\heb{#2})}

% Section markers
\newcommand{\peshat}{%
  \subsection*{\textcolor{peshatcolor}{On the Level of Peshat} \hfill {\hebrewfont\small\begin{hebrew}על דרך הפשט\end{hebrew}}}
}

\newcommand{\remez}{%
  \subsection*{\textcolor{remezcolor}{On the Level of Remez} \hfill {\hebrewfont\small\begin{hebrew}על דרך הרמז\end{hebrew}}}
}

\newcommand{\beeravraham}{%
  \begin{notebox}
  {\small\textsc{Be'er Avraham} — \textit{R' Avraham ben HaGRA}}\\[4pt]
}
\newcommand{\endbeeravraham}{\end{notebox}}

\newcommand{\radal}{%
  \begin{notebox}
  {\small\textsc{Pirush HaRaDaL} — \textit{R' David Luria}}\\[4pt]
}
\newcommand{\endradal}{\end{notebox}}

% Source reference
\newcommand{\src}[1]{{\small\textsuperscript{[\textit{#1}]}}}

% === Misc ===
\usepackage{microtype}
\usepackage{enumitem}
\setlist{nosep, leftmargin=*}

% Paragraph formatting
\setlength{\parindent}{0pt}
\setlength{\parskip}{6pt}

% === Hyperlinks ===
\usepackage[
  colorlinks=true,
  linkcolor=peshatcolor,
  urlcolor=peshatcolor
]{hyperref}

% ============================================================
\begin{document}

% === TITLE PAGE ===
\begin{titlepage}
\begin{center}
\vspace*{1.5in}

{\hebrewfont\Huge\begin{hebrew}פירוש הגר"א על מגילת אסתר\end{hebrew}}

\vspace{0.3in}

{\Huge\bfseries The Vilna Gaon\\[4pt] on Megillas Esther}

\vspace{0.4in}

{\large\itshape A Complete English Rendition}

\vspace{0.15in}

{\normalsize Including the commentary\\
\hebterm{Al Derekh HaPeshat}{על דרך הפשט} — On the Level of Plain Meaning\\
\hebterm{Al Derekh HaRemez}{על דרך הרמז} — On the Level of Allusion}

\vspace{0.4in}

{\small With notes from\\[4pt]
\textsc{Be'er Avraham} — R' Avraham son of the GRA\\
\textsc{Pirush HaRaDaL} — R' David Luria}

\vspace{\fill}

{\footnotesize Based on the Warsaw Edition\\
Source: HebrewBooks.org \#39741\\[8pt]
Digitized and rendered for Purim 5786}

\end{center}
\end{titlepage}

% === FRONT MATTER ===
\frontmatter

\chapter*{Preface}

This volume presents the complete commentary of the Vilna Gaon (Rabbi Eliyahu of Vilna, 1720--1797) on \hebterm{Megillas Esther}{מגילת אסתר}, rendered into English for the first time in its entirety.

The Gaon's commentary operates on two distinct levels:

\begin{itemize}
\item \textbf{Al Derekh HaPeshat} (\heb{על דרך הפשט}) — the straightforward interpretation, dealing with the historical narrative, the political dynamics, and the revealed miracles of the Purim story.

\item \textbf{Al Derekh HaRemez} (\heb{על דרך הרמז}) — the allusive interpretation, in which the Gaon reads the Megillah as an allegory for the soul's battle against the \hebterm{yetzer hara}{יצר הרע}, the structure of prayer, and the hidden spiritual architecture of redemption.
\end{itemize}

The remez layer is one of the Gaon's most remarkable achievements in biblical commentary. He maps the entire Megillah onto the daily order of prayer (\hebterm{seder ha-tefillah}{סדר התפילה}), revealing how each event in the Purim narrative corresponds to a stage in the Jew's daily spiritual ascent.

Also included are notes from R' Avraham ben HaGRA (\textsc{Be'er Avraham}), who supplements his father's terse comments with expanded explanations, and from R' David Luria (\textsc{Pirush HaRaDaL}), who provides cross-references and clarifications.

\vspace{12pt}
\begin{flushright}
\textit{Purim 5786 / March 2026}
\end{flushright}

\tableofcontents

% === MAIN MATTER ===
\mainmatter

% Chapters will be \input here
\input{chapter01}
% \input{chapter02}
% \input{chapter03}
% \input{chapter04}
% \input{chapter05}
% \input{chapter06}
% \input{chapter07}
% \input{chapter08}
% \input{chapter09}
% \input{chapter10}

\end{document}
